\section{Raft and XATMI Transaction Thinking}
\label{sec:RaftAndXatmi}

This section is for application developers and architects familiar with
XATMI, XA, two-phase commit, service transactions, or similar middleware
transaction models. Raft and XATMI can look similar from a distance because
both coordinate distributed work, but they are not the same mechanism.

\subsection{Short version}

Raft answers: \emph{Which command is committed at log index N, in what
order, according to a quorum?}

XATMI/XA answers: \emph{Can this transaction commit atomically across a
set of enlisted services?}

Raft is a replicated state-machine consensus mechanism. XATMI is an
application service and transaction coordination model. In this project, a
Raft application command is closer to a durable, replicated service
request than to a full XA transaction.

\subsection{Where they are similar}

Both models have a request that crosses process boundaries
(\texttt{client $\to$ service/runtime $\to$ coordinated outcome $\to$
response}) and both must handle timeouts, retries, process crashes,
network partitions, durable recovery, duplicate delivery, unclear client
outcome, and authentication/authorization before mutation.

Both force the application to answer the same practical question: if the
client did not receive the response, did the operation happen?

\subsection{Where they differ}

\subsubsection{Raft commits log entries}

Raft takes an application command and appends it to a replicated log.
Once the leader has replicated the entry to a quorum, the entry is
committed. Every correct replica applies committed entries in the same
order. The replicated log is the serialization point.

\subsubsection{XATMI/XA coordinates resources}

XATMI-style transaction processing coordinates a unit of work across
services and resource managers. XA/two-phase commit asks all enlisted
resources to prepare and then commit or roll back. The transaction
coordinator tries to make multiple independent participants agree on one
atomic outcome.

\subsubsection{The core difference}

Raft: one replicated log orders commands for one replicated state machine.

XATMI/XA: one transaction coordinates multiple participants/resources.

Raft gives you linearizable replication of a state machine. It does not
automatically give you an ACID transaction across arbitrary external
systems.

\subsection{Commit semantics}

\subsubsection{Raft commit}

A Raft command is committed when durably replicated to a quorum. After
commit, replicas apply it in log order. The application does not decide
whether an entry is committed---Raft does. \texttt{apply()} is called
only for committed log entries.

\subsubsection{XA commit}

An XA transaction commits when all resources have successfully passed the
transaction protocol. If a resource cannot prepare or commit, the
transaction manager must resolve the outcome using XA recovery rules. A
Raft application state machine does not participate in prepare/commit/
rollback.

\subsection{Ambiguous outcomes}

Both models have ambiguous client outcomes. A client sends a request, the
leader commits the command, the leader crashes before the response reaches
the client, and the client sees a timeout. From the client's perspective,
the operation is uncertain.

The remedy is the same design discipline in both models: assign a stable
request id, store enough outcome information to recognize retries, and
make retry behavior idempotent from the client's perspective.

\subsection{Idempotency and request identity}

Application developers should treat request identity as part of the domain
command model. A good command shape includes a unique
\texttt{commandId}, the \texttt{requesterId}, the operation type, the
domain identifier, and the payload. The state machine keeps a small
outcome table mapping \texttt{commandId $\to$ result}.

On first application: if the commandId is unknown, apply the domain
mutation, store the result, and return it. On retry: if the commandId is
known, return the previously stored result. This pattern is useful in
both Raft and XATMI-style systems and is especially important when clients
retry after timeout.

\subsection{Determinism}

Raft application commands must be deterministic. Every node applies the
same committed command at the same log index; every node must produce
the same state transition. Avoid inside \texttt{apply()}: local
wall-clock decisions, random number generation, reading local files not
replicated by Raft, calling external services, checking local node
identity, or using non-deterministic iteration order.

Prefer: include timestamps in the command if the business event needs one,
include generated ids in the command, validate before submission, and
make \texttt{apply()} a pure transition over current replicated state
plus command bytes. XATMI services can also suffer from non-determinism,
but Raft makes it stricter because all replicas must replay the same
command stream.

\subsection{External resources}

The safest Raft application state machine does not call external systems
during \texttt{apply()}. If \texttt{apply()} calls a database, message
broker, REST service, or mainframe transaction service, the operation is
no longer just a replicated state-machine transition---you have introduced
another participant with its own failure and recovery semantics.

\subsubsection{Recommended pattern: Raft as source of truth}

For reference data and masterdata: a management system validates a change,
submits a command to Raft, Raft commits it, the state machine updates
local replicated state, and local readers use fast in-memory lookups.
The authoritative distributed value is the Raft log and state machine.
External stores can feed it or observe it, but should not be part of the
committed application transition unless deliberately modeled.

\subsubsection{Outbox pattern}

If a committed Raft command must trigger external side effects, use an
outbox approach: the \texttt{apply} updates replicated domain state and
records a pending side effect in replicated state; a separate worker
observes the pending side effect, performs the external action
idempotently, and submits a follow-up command marking the side effect
complete. This keeps Raft apply deterministic and gives side effects a
recoverable state.

\subsubsection{Avoid hidden XA inside Raft apply}

Do not casually start an XA transaction inside \texttt{apply()}, update
an external database, call a remote service, and commit the XA
transaction. That design combines Raft commit with a second transaction
protocol. It may be correct in a specific architecture, but it is no
longer a simple Raft state machine and needs a full failure analysis.

\subsection{Mapping concepts}

\begin{longtable}{p{48mm}p{52mm}p{52mm}}
\toprule
XATMI/XA concept & Raft-related concept & Important difference \\
\midrule
Service request & Client command/query &
Raft commands are replicated before application mutation. \\
\midrule
Transaction commit & Log entry commit &
Raft commit is quorum replication, not multi-resource prepare/commit. \\
\midrule
Resource manager & Application state machine &
State machine should be deterministic and local. \\
\midrule
Transaction id & Command/request id &
Application should model idempotency explicitly. \\
\midrule
Timeout & Timeout &
Both can leave client outcome ambiguous. \\
\midrule
Recovery log & Raft log and snapshots &
Raft log orders the replicated state machine. \\
\midrule
Rollback & Compensating commands &
Raft commands are not rolled back after commitment; compensating commands
are used. \\
\midrule
Service routing & Leader redirect / client retry &
Writes must reach the leader. \\
\midrule
Read-only service & Linearizable query &
Raft read safety must be established first. \\
\bottomrule
\end{longtable}

\subsection{Rollback and compensation}

In XA, rollback is part of the transaction protocol before commit. In
Raft, once an entry is committed and applied, the application does not
roll it back. If the business process needs reversal, submit a new
compensating command. That compensation is itself a new committed fact in
the log.

\subsection{Reads and lookup workloads}

XATMI applications often think in terms of service calls. In this Raft
design, the target use case is often few writes, many reads. For
reference data and masterdata distribution: central governance produces
low-rate replicated changes through Raft; high-rate local lookups are
served from domain state. A local lookup is safe only after Raft has
established that the node can serve the read. The runtime does that before
calling the application's query method. Application developers should not
implement quorum checks in the query code.

\subsection{Cluster configuration is not domain data}

Cluster configuration includes current voters, next voters during joint
consensus, learners, membership transition state, and quorum rules. That
state belongs to Raft. It is not reference data, not masterdata, not a
key/value entry owned by the application. The application snapshot stores
business state; the Raft snapshot wrapper stores cluster configuration.
Applications must not reserve keys such as \texttt{\_raft/config} or
\texttt{\_cluster/voters}---those would blur the safety boundary.

\subsection{Transaction boundaries for domain developers}

When designing a domain command, ask:

\begin{enumerate}
\item What is the smallest deterministic state change?
\item What command id identifies this request across retries?
\item What result should be returned if the same command id is seen again?
\item What state must be included in snapshots?
\item What external side effects are needed, if any?
\item Can side effects be modeled as follow-up commands?
\item What authorization/admission rule applies before the command enters
Raft?
\end{enumerate}

Do not ask the application to decide: who the leader is, what quorum
means, whether a membership change is committed, where cluster
configuration is stored, or whether a read lease is valid. Those are Raft
runtime questions.

\subsection{Design guidance}

Use Raft when: the application can be represented as a deterministic state
machine, the replicated state is small or compactable enough for
snapshots, writes are lower volume than reads, local read performance
matters, and the system benefits from clear leader/quorum semantics.

Use XA/XATMI-style coordination when: one business transaction must
atomically update multiple independent resource managers, participants
cannot be represented as one replicated state machine, and prepare/
commit/rollback across external systems is the actual requirement.

Use both only with explicit design: Raft for replicated state, XA/XATMI
for external systems that truly require transaction coordination,
idempotent bridges between them, and clear recovery and compensation
rules.

\subsection{Practical rule}

For a domain application developer using this project: your command is the
transaction-like business request; Raft makes that request durable,
ordered, replicated, and recoverable; your state machine makes the
deterministic business change; Raft cluster configuration stays outside
your domain data.
